% =====================================================================
%  VII. FUTURE SCOPE
%  The directions below follow from the design and its limitations
%  (Section VI-B). Directions that depend on the results are added, and
%  any of these that the results undercut are dropped, once Section V
%  exists.
% =====================================================================
\section{Future Scope}
\label{sec:future}

Several directions follow from the limitations in
Section~\ref{sec:limitations}.

\emph{Observing escalation.} The constructed policies and the swept
detection probabilities are the least grounded part of the comparator. A
study that records how people move between the configurations of a tiered
service, including where they start, how often they notice a wrong answer
and when they give up, would replace those assumptions with estimates. It
would turn the computed tax into a measured one and give P4 an observed
starting distribution.

\emph{Continuous reasoning budgets.} Qwen3 and several hosted services let
the amount of reasoning be set as a token budget or an effort level rather
than switched on or off. Treating the budget as a continuous axis would turn
the eight-cell grid into a surface, and the router's decision into a choice
of size and budget. The dependence of our labels on the thinking cap points
the same way, since the cap is a budget fixed once for every prompt.

\emph{Conversations and switching costs.} A cost model in which a retry
carries the conversation so far, and in which moving to another size pays
for loading it, would bring the tax closer to what users of memory-limited
hardware pay. It would also make the best choice for a query depend on the
model already loaded, which a one-shot router of the kind studied here does
not consider.

\emph{Routing the wider pools.} The public and hosted arms show what the tax
costs on pools that mix model families and frontier tiers, but both report
manual policies and the oracle rather than a trained router. Adding one
means recovering each public query's text from its source dataset and
training on outcomes we did not generate, or paying to label enough hosted
prompts to train on. Doing either, and adding quantisation as a third axis
beside size and mode, would test whether the minimum sufficient
configuration can still be predicted when cells differ in more than size and
mode.

\emph{Whole-system energy.} Measuring processor and memory power alongside
the GPU would close the energy boundary, and it would give CPU routers a
measured cost in place of an upper bound.

\emph{Online adaptation.} Providers change models and prices, and a router
trained on one pool's labels goes stale. The outcomes of routed queries,
as revealed by whether the user accepts an answer or retries, could keep a
router current without relabelling every prompt in every cell. That is a
partial-feedback problem, because only the routed cell's outcome is ever
observed.

\tbd{Add directions motivated by the results, and drop any of the above
that the results undercut, once Section~\ref{sec:results} exists.}
